\documentclass[10pt,a4paper]{amsart}
\usepackage[margin=19mm]{geometry}
\usepackage{amsmath,amssymb,mathtools}
\usepackage[scr=rsfs,cal=euler]{mathalfa}
\usepackage{xfrac}
\usepackage[unicode,hidelinks,bookmarks=false]{hyperref}
\newcommand{\ie}{i.e.\ }
\newcommand{\cf}{cf.\ }
\newcommand{\resp}{respectively\ }
\newcommand{\Subsection}{\subsection}
\providecommand{\nobreakdash}{\nobreak\hbox{-}}
\setlength{\emergencystretch}{2em}
\multlinegap0pt
\usepackage{bm}
\usepackage{bbm}
\usepackage{dsfont}
\usepackage{xspace}
\usepackage{cancel}
\usepackage{mathtools}
\usepackage{amsthm}
\makeatletter
\@ifundefined{theorem}{\newtheorem{theorem}{Theorem}}{}
\@ifundefined{lemma}{\newtheorem{lemma}[theorem]{Lemma}}{}
\makeatother
\title[Solutions to Three Conjectures and an Open Problem on Binary]{Solutions to Three Conjectures and an Open Problem on Binary BCH Codes}
\author{Xiaoqiang Wang, Jiawei He, Boru Yi, Dabin Zheng}
\date{Source: arXiv:2609.00532v2 (CC BY 4.0)}
\begin{document}
\maketitle

\noindent\emph{Source and licence.} Solutions to Three Conjectures and an Open Problem on Binary BCH Codes. Version arXiv:2609.00532v2 (CC BY 4.0), distributed under the Creative Commons Attribution 4.0 International licence, as recorded in the arXiv licence metadata for this exact version and shown on the abstract page https://arxiv.org/abs/2609.00532v2. Full rights notice: licenses/arxiv-2609.00532-CC-BY-4.0.txt.\par\smallskip

\noindent\textit{Editorial scope.} Counted unit: the Lemma whose source label is \texttt{lem:V1I} in the article (source0.tex lines 879--905, source label \texttt{lem:V1I}), delivered together with the article's own complete original proof of it (source0.tex lines 907--1361, one \texttt{proof} environment of 455 lines). No mathematics is written, repaired or re-derived: statement and proof are the article's own text line for line. Required context retained uncounted: none. Every definition, display and label of those lines is retained unchanged. Omitted from this package and not redistributed: 1--878, 906--906, 1362--2181. Nothing inside a retained excerpt is omitted or altered: every retained body is delivered exactly as the article's lines are written, so no omission receipt of any kind accompanies this package. Citations: the retained excerpts carry no citation command at all, so no bibliography entry of the article is transcribed and no reference is redistributed. Reference adaptation: none was needed and none was made; every reference inside the delivered unit resolves inside it, so no unresolved reference or citation is produced. Printed slips kept literal: no printed word, symbol, hypothesis or number of the counted statement or of its proof was altered. Layout and class: the article's own class is \texttt{IEEEtran} with packages including amssymb, amsthm, amsmath, latexsym, graphicx, mathrsfs, amsfonts, longtable, setspace, caption, rotating, ifsym, bbm, makecell; the wrapper is the lane's pinned \texttt{amsart} editorial wrapper at 10pt on A4, which restates the theorem-like environments the retained text needs and adds the article's own macro definitions that the retained text needs (none), with extra wrapper packages (bm, bbm, dsfont, xspace, cancel, mathtools). The delivered numbering is the unit's own. Assets: no retained span contains a figure environment, a graphics-inclusion command, a PDF-page inclusion, a table or a TikZ picture; no image file is copied into this package, the package declares no asset dependency and no image file is redistributed with it. Licence: version arXiv:2609.00532v2 of this article is distributed under the Creative Commons Attribution 4.0 International licence, as recorded in the arXiv licence metadata for this exact version and shown on the abstract page https://arxiv.org/abs/2609.00532v2. A full rights notice is delivered at licenses/arxiv-2609.00532-CC-BY-4.0.txt. Characters: every retained excerpt is pure ASCII, so the delivered engine, XeLaTeX with its default Latin Modern text font, reports no missing character.
\begin{lemma}\label{lem:V1I}
Let \(s\) be a positive integer such that
\[
s\equiv 2 \pmod 6,
\quad\text{or}\quad
s\equiv 4 \pmod 6,
\quad\text{or}\quad 5\mid s
\,\,\text{and}\,\,
15\nmid s
\]
and set
$
n=2^{2s}+2^s+1.
$
Then there exist five pairwise distinct elements
$
x,y,z,u,v\in \mathbb F_{2^{3s}}
$
such that
\begin{eqnarray*}
\begin{cases}
x+y+z+u+v=0,  \\
x^3+y^3+z^3+u^3+v^3=0,\\
x^n=y^n=z^n=u^n=v^n=1.
\end{cases}
\end{eqnarray*}
\end{lemma}

\begin{proof}
We prove this result from the following two cases.

\noindent
\textbf{Case 1:} $s\equiv 2 \pmod 6$,
or
$s\equiv 4 \pmod 6$.  We first show that \(21\mid n\).
Since \(s\) is even, we have
$
2^s\equiv 1 \pmod 3,
$
and hence
$
2^{2s}\equiv 1 \pmod 3.
$
Therefore,
\[
n=2^{2s}+2^s+1\equiv 1+1+1\equiv 0 \pmod 3,
\]
which means that \(3\mid n\).
Now consider the residue of \(2^s\) modulo \(7\). The multiplicative order of \(2\) modulo \(7\) is \(3\). Since \(s\equiv 2\) or \(4\pmod 6\), we have \(3\nmid s\), and thus \(2^s\not\equiv 1\pmod 7\). Moreover, \((2^s)^3\equiv 1\pmod 7\), so \(2^s\) is a nontrivial cube root of unity modulo \(7\). It follows that
\[
(2^s)^2+2^s+1\equiv 0 \pmod 7.
\]
Hence, \(7\mid n\). It follows that
$
21\mid n.
$

Consider the polynomial
\begin{equation}\label{eqhX}
h(X)=X^6+X^4+X^2+X+1\in \mathbb F_2[X].
\end{equation}
Since \(\deg {h(X)}=6\), if \(h(X)\) were reducible over \(\mathbb F_2\), then it would have an irreducible factor of degree at most \(3\). We therefore exclude irreducible factors of degrees \(1\), \(2\), and \(3\).

First,
$h(0)=1$ and $h(1)=1,$
so \(h(X)\) has no linear factor over \(\mathbb F_2\).
The unique monic irreducible quadratic polynomial over \(\mathbb F_2\) is
$$
f_2(X)=X^2+X+1.
$$
It is clear that
$
X^2\equiv X+1 \pmod {f_2(X)}
$ and $
X^3\equiv1\pmod {f_2(X)}.
$
Then
$
X^4\equiv X \pmod {f_2(X)}
$ and
$X^6\equiv1 \pmod {f_2(X)}.
$
Hence,
$$
h(X)
\equiv
1+X+(X+1)+X+1
=
X+1
\not\equiv0
\pmod{f_2(X)}.
$$
Therefore, \(f_2(X)\nmid h(X)\).
There are exactly two monic irreducible cubic polynomials over \(\mathbb F_2\):
$$
f_{3,1}(X)=X^3+X+1\,\, \text{and}\,\,
f_{3,2}(X)=X^3+X^2+1.
$$
It ia obvious that
$
X^3\equiv X+1 \pmod {f_{3,1}(X)}.
$
Thus,
$
X^4\equiv X^2+X \pmod {f_{3,1}(X)},
$ and $
X^6=(X^3)^2\equiv X^2+1 \pmod {f_{3,1}(X)}.
$
It follows that
$$
\begin{aligned}
h(X)
&\equiv
(X^2+1)+(X^2+X)+X^2+X+1\\
&=X^2
\not\equiv0
\pmod{f_{3,1}(X)}.
\end{aligned}
$$
Hence, \(f_{3,1}(X)\nmid h(X)\).

Similarly, we have
$
X^3\equiv X^2+1 \pmod {f_{3,2}(X)}.
$
Then
$
X^4\equiv X^2+X+1 \pmod {f_{3,2}(X)}
$
and
$
X^6\equiv X^2+X \pmod {f_{3,2}(X)}.
$
Hence,
$$
\begin{aligned}
h(X)
&\equiv
(X^2+X)+(X^2+X+1)+X^2+X+1\\
&=X^2+X
\not\equiv0
\pmod{f_{3,2}(X)}.
\end{aligned}
$$
Thus, \(f_{3,2}(X)\nmid h(X)\). Hence, \(h(X)\) has no irreducible factor of degree \(1\), \(2\), or \(3\). Since \(\deg h(X)=6\), every nontrivial factorization of \(h(X)\) would contain a factor of degree at most \(3\). Therefore, (\ref{eqhX})
is irreducible over \(\mathbb F_2\).


Now let \(\theta\) be a root of \(h(X)\). Since
$
h(\theta)=\theta^6+\theta^4+\theta^2+\theta+1=0,
$
we have
\begin{equation}\label{eq:0831}
\theta^6=\theta^4+\theta^2+\theta+1.
\end{equation}
Multiplying both sides by \(\theta\), we obtain
$
\theta^7=\theta^5+\theta^3+\theta^2+\theta.
$
Squaring this identity in characteristic \(2\) gives
\begin{equation}\label{eq:083101}
\theta^{14}
=
\theta^{10}+\theta^6+\theta^4+\theta^2.
\end{equation}
Combining (\ref{eq:0831}) and (\ref{eq:083101}),
we further obtain
$
\theta^8
=
\theta^3+\theta+1,
$
and hence
$
\theta^{10}
=
\theta^5+\theta^3+\theta^2.
$
Substituting these expressions into the formula for \(\theta^{14}\), we get
$$
\begin{aligned}
\theta^{14}
&=
\bigl(\theta^5+\theta^3+\theta^2\bigr)
+\bigl(\theta^4+\theta^2+\theta+1\bigr)
+\theta^4+\theta^2=
\theta^5+\theta^3+\theta^2+\theta+1=
\theta^7+1.
\end{aligned}
$$
Then
$$
\theta^{21}
=
\theta^{14}\theta^7
=
(\theta^7+1)\theta^7
=
\theta^{14}+\theta^7
=
(\theta^7+1)+\theta^7
=
1.
$$
Thus, the multiplicative order of \(\theta\) divides \(21\).

Since \(h(X)\) is irreducible over \(\mathbb F_2\) and has degree \(6\), the minimal polynomial of \(\theta\) over \(\mathbb F_2\) has degree \(6\). Hence, the multiplicative order of \(\theta\) cannot be \(1\), \(3\), or \(7\). Indeed, if \(\theta^3=1\), then \(\theta\) would lie in \(\mathbb F_{2^2}\), so its minimal polynomial over \(\mathbb F_2\) would have degree at most \(2\). Similarly, if \(\theta^7=1\), then \(\theta\) would lie in \(\mathbb F_{2^3}\), so its minimal polynomial would have degree at most \(3\). Both cases contradict the fact that the minimal polynomial of \(\theta\) has degree \(6\).
Therefore, the multiplicative order of \(\theta\) is exactly
$
21.
$
Since every root of the irreducible polynomial \(h(X)\) is a Frobenius conjugate of \(\theta\), all roots of \(h(X)\) have the same multiplicative order. Hence every root of \(h(X)\) has multiplicative order \(21\).

Set
\[
x=1,\quad
y=\theta,\quad
z=\theta^6,\quad
u=\theta^8,\quad
v=\theta^{18}.
\]
The exponents \(0,1,6,8,18\) are pairwise distinct modulo \(21\), so the five elements \(x,y,z,u,v\) are pairwise distinct. Since \(21\mid n\), we have \[
x^n=y^n=z^n=u^n=v^n=1.
\]
It remains to verify the two additive conditions. A direct computation in \(\mathbb F_2[X]\) gives
\[
h(X)\mid 1+X+X^6+X^8+X^{18}.
\]
Evaluating at \(X=\theta\), we obtain
$
1+\theta+\theta^6+\theta^8+\theta^{18}=0,
$
that is,
\[
x+y+z+u+v=0.
\]
Similarly,
\[
\begin{aligned}
x^3+y^3+z^3+u^3+v^3
=1+\theta^3+\theta^{18}+\theta^{24}+\theta^{54}=1+\theta^3+\theta^{18}+\theta^3+\theta^{12}=1+\theta^{12}+\theta^{18},
\end{aligned}
\]
where we used \(\theta^{21}=1\) and the fact that the underlying field has characteristic \(2\). Another direct computation shows that
$
h(X)\mid 1+X^{12}+X^{18}.
$
Then
$
1+\theta^{12}+\theta^{18}=0,
$
and hence
\[
x^3+y^3+z^3+u^3+v^3=0.
\]

\noindent
\textbf{Case 2:} $5\mid s$ and $15\nmid s$. Write
$
s=5k.
$
Since \(15\nmid s\), we have \(3\nmid k\). Put
$
M=2^{10}+2^5+1=1057.
$
Since
\[
2^{15}-1=(2^5-1)(2^{10}+2^5+1)=31M,
\]
we have
$
M\mid 2^{15}-1.
$
We first show that \(M\mid n\). Since \(3\nmid k\), we have
\[
k\equiv 1 \pmod 3
\quad\text{or}\quad
k\equiv 2 \pmod 3.
\]
If \(k\equiv 1\pmod 3\), then
$
2^{5k}\equiv 2^5\pmod M
$
and
$
2^{10k}\equiv 2^{10}\pmod M.
$
Hence,
\[
n=2^{10k}+2^{5k}+1
\equiv 2^{10}+2^5+1
\equiv 0\pmod M.
\]
If \(k\equiv 2\pmod 3\), then
$
2^{5k}\equiv 2^{10}\pmod M
$
and
$
2^{10k}\equiv 2^{20}\equiv 2^5\pmod M.
$
Thus again
\[
n\equiv 2^5+2^{10}+1\equiv 0\pmod M.
\]
Therefore,
$
M\mid n.
$

Since \(M\mid 2^{15}-1\), there exists an element
$
\theta\in \mathbb F_{2^{15}}^{*}
$
of multiplicative order \(M=1057\). As \(5\mid s\), we have
$
15\mid 3s,
$
and hence
$
\mathbb F_{2^{15}}\subseteq \mathbb F_{2^{3s}}.
$
Thus \(\theta\in\mathbb F_{2^{3s}}\).
Let
\[
m(X)
=
X^{15}+X^{10}+X^9+X^8+X^4+X^3+X^2+X+1
\in\mathbb F_2[X].
\]
We prove that \(m(X)\) is irreducible over \(\mathbb F_2\).
Assume \(m(X)=0\). We will prove that \(X\) has order \(1057\) in \(\mathbb F_{2^{15}}^\times\). Consequently,
\[
X^{15}=X^{10}+X^9+X^8+X^4+X^3+X^2+X+1.
\]
Hence,
\[
X^{16}=X^{11}+X^{10}+X^9+X^5+X^4+X^3+X^2+X.
\]
Over \(\mathbb F_2\), it is easy to check that
\[
X^{32}=X^{22}+X^{20}+X^{18}+X^{10}+X^8+X^6+X^4+X^2.
\]
Successive reduction modulo \(m(X)\) gives
\[
X^{32}\equiv X^{14}+X^9+X^8+X^7+X^3+X^2+X\pmod {m(X)}.
\]
Thus,
\[
X^{33}\equiv X^{15}+X^{10}+X^9+X^8+X^4+X^3+X^2\pmod {m(X)}.
\]
Substituting the reduction of \(X^{15}\) and cancelling equal terms in characteristic \(2\),
\[
X^{33}\equiv X+1 \pmod {m(X)}.
\]
Since \(1056=32\cdot 33\), we have
\[
X^{1056}=(X^{33})^{32}\equiv (X+1)^{32}=X^{32}+1\pmod {m(X)}.
\]
Multiplying by \(X\),
\[
X^{1057}\equiv X^{33}+X\equiv (X+1)+X=1\pmod {m(X)}.
\]
Therefore,
$
m(X)\mid X^{1057}-1.
$

Since
$
2^{15}-1=32767=31\cdot 1057,
$
we have \(1057\mid 2^{15}-1\), then \(X^{1057}-1\mid X^{2^{15}-1}-1\). Hence,
\[
m(X)\mid X^{2^{15}}-X.
\]
The polynomial \(X^{2^{15}}-X\) is squarefree over \(\mathbb F_2\), and its irreducible factors are precisely the monic irreducibles over \(\mathbb F_2\) whose degrees divide \(15\). Therefore, \(m(X)\) is squarefree and every irreducible factor of \(m(X)\) has degree \(d\mid 15\). The positive divisors of \(15\) are \(1,3,5,15\); it remains to exclude \(1,3,5\).
Since
\[
m(0)=1\neq 0\,\, \text{and}\,\, m(1)=9\equiv 1\not\equiv 0\pmod 2,
\]
\(m(X)\) has no linear factor over \(\mathbb F_2\).
A cubic irreducible factor would divide \(X^8-X\). It is easy to check that
\[
m(X)\equiv X^4+X+1\pmod{X^8-X},
\]
and the Euclidean algorithm gives
\[
\gcd(X^8-X,X^4+X+1)=1.
\]
Thus \(m(X)\) has no cubic irreducible factor.
A quintic irreducible factor would divide \(X^{32}-X\). Let
\[
R(X)=X^{14}+X^9+X^8+X^7+X^3+X^2
\]
be the remainder of \(X^{32}+X\) modulo \(m(X)\). Since \(X^{32}-X=X^{32}+X\) in characteristic \(2\),
\[
\gcd(m(X),X^{32}-X)=\gcd(m(X),R(X)).
\]
Moreover,
$
m(X)+XR(X)=X^2+X+1,
$
then
\[
\gcd(m(X),R(X))=\gcd(R(X),X^2+X+1).
\]
Substitution gives \(R(X)\equiv 1\pmod{X^2+X+1}\), we have
$
\gcd(m(X),X^{32}-X)=1.
$
Thus, \(m(X)\) has no quintic irreducible factor.
Consequently, every irreducible factor of \(m(X)\) has degree \(15\). Since \(\deg m(X)=15\), \(m(X)\) is irreducible over \(\mathbb F_2\), and
\[
\mathbb F_2[X]/(m(X))\cong \mathbb F_{2^{15}}.
\]
It remains to determine the order of \(X\) in \(\mathbb F_{2^{15}}^{*}\). Since \(1057=7\cdot 151\) with \(7\) and \(151\) prime, and since \(X^{1057}\equiv 1\pmod {m(X)}\), the order of \(X\) is \(1057\) if and only if
\[
X^7\not\equiv 1\pmod {m(X)}\,\,\,\, \text{and}\,\,\,\, X^{151}\not\equiv 1\pmod {m(X)}.
\]
The first is immediate: \(X^7-1\) has degree \(7\), so it cannot be divisible by the degree-\(15\) polynomial \(m(X)\).
For the second, repeated squaring yields
\[
X^{64}\equiv 1+X^3+X^6+X^9+X^{13}+X^{14}\pmod {m(X)},
\]
\[
X^{128}\equiv 1+X+X^4+X^7+X^{11}+X^{12}\pmod {m(X)},
\]
and, using
\[
X^{23}\equiv X+X^3+X^4+X^5+X^7+X^8+X^{10}+X^{12}+X^{13}\pmod {m(X)},
\]
one obtains
\[
X^{151}=X^{128}X^{23}\equiv X^{14}+X^{10}+X^6+X^3+X \pmod {m(X)}.
\]
This is not the constant polynomial \(1\), so \(X^{151}\not\equiv 1\pmod {m(X)}\). Therefore the order of \(X\) in \(\mathbb F_{2^{15}}^\times\) is exactly \(1057\).
Thus
we may choose \(\theta\) such that \(m(\theta)=0\) and
\(\operatorname{ord}(\theta)=1057\). A direct computation in
\(\mathbb F_2[X]\) gives
$$
m(X)\mid
\left(
1+X^{72}+X^{105}+X^{119}+X^{234}
\right)
\,\,\text{and}\,\,
m(X)\mid
\left(
1+X^{216}+X^{315}+X^{357}+X^{702}
\right).
$$
Therefore,
$$
1+\theta^{72}+\theta^{105}+\theta^{119}+\theta^{234}=0
\,\,\text{and}\,\,
1+\theta^{216}+\theta^{315}+\theta^{357}+\theta^{702}=0.
$$

Now set
\[
x=1,\qquad
y=\theta^{72},\qquad
z=\theta^{105},\qquad
u=\theta^{119},\qquad
v=\theta^{234}.
\]
We have
\[
x+y+z+u+v=0
\,\,\text{and}\,\,
x^3+y^3+z^3+u^3+v^3=0.
\]

Finally, since each of \(x,y,z,u,v\) has multiplicative order
dividing \(M\), and \(M\mid n\), we obtain
\[
x^n=y^n=z^n=u^n=v^n=1.
\]

From the above two cases, we obtain the desired result.
\end{proof}

\end{document}
